\documentclass{plantillaPPT}

\titulo{2}{Límites y Continuidad}

\begin{document}
\primeraDiapo

\begin{frame}{Límites}
\framesubtitle{Práctica 2}

1. Sea \(f\) la siguiente función:
\[f(x,y)=\frac{x-2y}{4y^2-x^2}, \quad 4y^2\neq x^2\]

Calcule, si existe, \(\displaystyle \lim_{(x,y)\to(a,b)}\) cuando:\\[10 pt]
\((a) \quad (a,b)=(2,1)\)\\
\((b) \quad (a,b)=(0,0)\)\\
\((c) \quad (a,b)=(2,-1)\)

\end{frame}

\begin{frame}{Continuidad}

\begin{definicion}
Sea \(f:\mathbb{R}^n \to \mathbb{R}\) y un punto \(\Vec{x_0}\in\mathbb{R}^n\).\\
Se dice que \(f\) es continua en \(\vec{x_0}\) si:
\[
\lim_{\vec{x}\to\vec{x_0}} f(\vec{x}) = f(\vec{x_0})
\]
\end{definicion}
    
\end{frame}

\begin{frame}{Resultado}
\framesubtitle{Continuidad}

De la definición se obtienen 3 condiciones:

\begin{enumerate}
    \item \(f(\vec{x_0})\) existe.
    \vspace{5pt}
    \item \(\displaystyle \lim_{\vec{x}\to\vec{x_0}}f(\vec{x})\) existe.
    \vspace{5pt}
    \item \(\lim_{\vec{x}\to\vec{x_0}} f(\vec{x}) = f(\vec{x_0})\).
\end{enumerate}

\end{frame}

\begin{frame}{Límites en un punto}
\framesubtitle{Definición formal para campos escalares}

\begin{definicion}
    Sea \(f:D\subseteq\mathbb{R}^n\to\mathbb{R}\) una función escalar y \(\vec{x_0}\) un punto de acumulación de \(D\). Se dice que:
    \[
    \lim_{\vec{x}\to\vec{x_0}}f(\vec{x}) = L
    \]
    si se cumple que
    \[\forall \epsilon>0, \exists \delta>0: \|\vec{x}-\vec{x_0}\|<\delta \Rightarrow \|f(\vec{x})-L\|<\epsilon\]
\end{definicion}

\end{frame}

\begin{frame}{Definición de límite}
\framesubtitle{Práctica 2}

2. Use la definición de límite para mostrar que \(f\) es continua en \((0,0)\), donde:
\[
f(x) =
\begin{cases}
\frac{x^2-y\cos(y)}{|x|+\sqrt{|y|}}, \quad \;\text{si }(x,y)\neq(0,0) \\
 \quad \quad 0, \quad \quad \quad \text{si }(x,y) = (0,0)
\end{cases}
\]
    
\end{frame}

\begin{frame}{Continuidad}
\framesubtitle{Práctica 2}

3. Se da una función \(f(x,y)\) que no está definida en \((0,0)\). ¿Es posible definir el valor \(f(0,0)\) de tal modo que \(f\) sea continua en dicho punto?. Justifique su respuesta.
\[(a) f(x,y)=\frac{xy^2}{x^2+y^4}\]
\[(b) f(x,y)=\frac{\sin(x^2+y^2)}{2(e^{x^2+y^2}-1)}\]
    
\end{frame}

\begin{frame}{Continuidad}
\framesubtitle{Práctica 2}

4. Sea \(f: \mathbb{R}^2\to\mathbb{R}\) una función definida por:
\[
f(x) = 
\begin{cases}
\frac{x^2+xy}{x-y}, \quad si x\neq y \\
\quad 0, \quad \quad si x=y 
\end{cases}
\]
\((a)\) Determinar el subconjunto mayor de \(\mathbb{R}^2\) en donde \(f\) es continua.\\
\end{frame}

\end{document}